\section{Quantum algorithms}

We now review standard Quantum algorithms in the tensor network formalism.

\subsection{Fourier transform}

Recall the tensor network decomposition of the Fourier transform developed in \secref{sec:fourierTransform}.

When $\vectorat{\shortcatvariables}$ is a quantum state, the Fourier transform can be understood as a unitary evolving the state.
Since both the controlled $U$ tensor and the Hadamard $H$ are unitary, each of them can be implemented by a sequence of gates.
For sparse implementation one can exploit that $U$ is a tensor product of matrices (i.e. single qubit gates).

% Gate complexity
To realize the QFT we need $\catorder$ single-qubit Hadamard gates and $\sum_{\catenumeratorin} \catenumerator = \frac{\catorder\cdot(\catorder-1)}{2}$ controlled diagonal single-qubit gates (to realize the $U^{\catenumerator}$ transforms).
The gate complexity is therefore $\mathcal{O}(\catorder^2)$.



\subsubsection{Detecting periods}

Based on the support of the Fourier transform, one can derive the periods of a function.
Periods are indices $(\thirdcatvariableof{0},\ldots,\thirdcatvariableof{\catorder-1})$ such that
\begin{align*}
    \hypercoreat{\shortcatvariables=\shortcatindices}
    = \hypercoreat{\shortcatvariables=\shortcatindices\oplus\thirdcatindexof{[\catorder]}}
\end{align*}
where $\oplus$ is the tuple position wise sum $\mathrm{mod}\,\,\catdimof{\catenumerator}$.


Each supported index of the Fourier transform corresponds with periods in the function.
The periods of the functions are thus at least the intersection of those (when assuming more structure one might have more periods).

\subsubsection{Walsh-Hadamard transform}

The Walsh-Hadamard transform is the Fourier transform with respect to the group $\bigtimes_{\catenumeratorin}[2]$.
It is the legwise application of the Hadamard matrix.

In this case, we can easily determine the periods corresponding to each basis vector.
The inverse Fourier transform of $\onehotmapofat{\selindexof{[\catorder]}}{\selvariableof{[\catorder]}}$ is the tensor
\begin{align*}
    \left(\bigotimes_{\catenumeratorin \wcols \selindexof{\catenumerator}=0} \hplusbasat{\catvariableof{\catenumerator}} \right)
    \otimes
    \left(\bigotimes_{\catenumeratorin \wcols \selindexof{\catenumerator}=1} \hminusbasat{\catvariableof{\catenumerator}} \right) \, .
\end{align*}

For each $\hplusbasat{\catvariableof{\catenumerator}}$ we have a period generator of $\onehotmapof{\catenumerator}$.
For each pair of $\hminusbasat{\catvariableof{\catenumerator}}\otimes \hminusbasat{\catvariableof{\seccatenumerator}}$ we have a period generator of $\onehotmapof{\catenumerator}+\onehotmapof{\seccatenumerator}$.

One example is the Deutsch-Jozsa test, where the transformed is the $\onehotmapofat{0}{\selvariableof{[\catorder]}}$ state, if and only if the function is constant.
In this case, any vector is a period, since the group itself is generated by the period generators.

\subsubsection{Application}

QFT is a typical component in Quantum algorithms: QFT can be done efficiently with $\mathcal{0}(\catorder^2)$ controlled single-qubit gates.
If the Fourier transform is a basis vector, a single measurement is sufficient to determine it.




\subsection{Quantum Phase Estimation}\label{sec:quantumPhaseEstimation}\glsadd{qc:QPE}

% Generic Quantum Phase Estimation
When having a circuit implementing a unitary, the phase of the eigenvalues can be estimated using the Quantum Phase Estimation routine \cite{kitaev_quantum_1995, cleve_quantum_1998}.

Let $\unitaryat{\ccatvariableof{\insymbol}{[\catorder]},\ccatvariableof{\outsymbol}{[\catorder]}}$ be a unitary then the absolute of any eigenvalue is $1$ and the eigenvalue is expressible by
\begin{align*}
    \expof{2\pi i \cdot \phi}
\end{align*}
for some $\phi\in[0,1]$.
Let $\phi=0.\phi_{0}\cdots\phi_{n-1}$ be a $2$-adic representation of $\phi$, that is
\begin{align*}
    \phi
    = \sum_{\catenumerator\in[n]} 2^{-\catenumerator-1} \phi_{\catenumerator} \, .
\end{align*}

In the Quantum Phase Estimation routine the Fourier transform of $\onehotmapofat{\phi_0,...,\phi_{n-1}}{\selvariableof{[n]}}$
\begin{align*}
    \sqrt{\frac{1}{2^{n}}} \cdot \bigotimes_{\catenumerator\in[n]} \left(\fbasisat{\selvariableof{0}} + \expof{2\pi i \cdot 0.\phi_{n-\catenumerator-1}\cdots\phi_{n-1}}\right)
\end{align*}
is prepared as explained below.
The inverse Fourier transform on the register $\selvariableof{[n]}$ results then in the state
\begin{align*}
    \onehotmapofat{\phi_0,...,\phi_{n-1}}{\selvariableof{[n]}} \otimes \qstateat{\shortcatvariables}
\end{align*}
and $\phi_{\catenumerator}$ for $\catenumerator\in[n]$ can be retrieved by a single computational basis measurement of the register $\selvariableof{[n]}$.

To prepare the Fourier transform of $\onehotmapofat{\phi_0,...,\phi_{n-1}}{\shortcatvariables}$, we start with the state
\begin{align*}
    \hplusbasat{\selvariable} \otimes \qstateat{\shortcatvariables} \, .
\end{align*}
When applying $\unitaryat{\ccatvariableof{\insymbol}{[\catorder]},\ccatvariableof{\outsymbol}{[\catorder]}}$ $2^{\catenumerator}$ times controlled on $\selvariable$ we get the state
\begin{align*}
    \sqrt{\frac{1}{2}} \left(\fbasisat{\selvariableof{0}} + \expof{2\pi i \cdot 0.\phi_{n-\catenumerator-1}\cdots\phi_{n-1}}\right) \otimes \qstateat{\shortcatvariables}
\end{align*}

Now we apply for each $\catenumerator\in[n]$ the unitary $\unitaryat{\ccatvariableof{\insymbol}{[\catorder]},\ccatvariableof{\outsymbol}{[\catorder]}}$ $2^{\catenumerator}$ times controlled on $\selvariableof{\catenumerator}$ and get
\begin{align*}
    \sqrt{\frac{1}{2^{n}}} \onesat{\selvariableof{[n]}} \otimes \qstateat{\shortcatvariables} \, .
\end{align*}

\subsection{Grover operator}\label{sec:groverOperator}

We derive quantum inference schemes based on the Grover operator, which we want to first introduce in full generality.

\subsubsection{Definition}

Let us assume we have two orthogonal quantum states $\qstateofat{\goodstatesymbol}{\shortcatvariables},\qstateofat{\badstatesymbol}{\shortcatvariables}$ spanning a subspace $\subspaceof{}=\spanof{\qstateofat{\goodstatesymbol}{\shortcatvariables},\qstateofat{\goodstatesymbol}{\shortcatvariables}}$ and let us choose another state
\begin{align*}
    \qstateofat{0}{\shortcatvariables}
    \in \subspaceof{}\, .
\end{align*}
Then there is $\rotanglesymbol\in[0,4\pi)$ with
\begin{align*}
    \qstateofat{0}{\shortcatvariables}
    = \sinof{\frac{\rotanglesymbol}{2}}\qstateofat{\goodstatesymbol}{\shortcatvariables}
    + \cosof{\frac{\rotanglesymbol}{2}}\qstateofat{\badstatesymbol}{\shortcatvariables} \, .
\end{align*}
Consider the reflection about the states $\qstateof{\badstatesymbol}$ and $\qstateof{0}$ by
\begin{align*}
    \reflectionofat{\badstatesymbol}{\inshortcatvariables,\outshortcatvariables}
    &= 2 \cdot \qstateofat{\badstatesymbol}{\inshortcatvariables} \otimes \qstateofat{\badstatesymbol}{\outshortcatvariables} - \identityat{\inshortcatvariables,\outshortcatvariables} \\
    \reflectionofat{0}{\inshortcatvariables,\outshortcatvariables}
    &= 2 \cdot \qstateofat{0}{\inshortcatvariables} \otimes \qstateofat{0}{\outshortcatvariables} - \identityat{\inshortcatvariables,\outshortcatvariables} \, ,
\end{align*}
and their concatenation
\begin{align*}
    \groverat{\inshortcatvariables,\outshortcatvariables}
    = \contractionof{
        \reflectionofat{\badstatesymbol}{\inshortcatvariables,\ccatvariableof{[\catorder]}{\midsymbol}},
        \reflectionofat{0}{\ccatvariableof{[\catorder]}{\midsymbol},\outshortcatvariables}
    }{\inshortcatvariables,\outshortcatvariables}
\end{align*}
which is called the \emph{Grover operator}.

\subsection{Action on projected quantum states}

The Grover operator leaves any quantum state $\qstate$ invariant, which is orthogonal to $\qstateof{\goodstatesymbol}$ and $\qstateof{\badstatesymbol}$ (such states get a factor $(-1)$ by both reflection operations).
The projection of the quantum state onto $\subspaceof{}$ gets rotated by the Grover operator (see \figref{fig:rotationSketch}).
To be more precise, we denote the projection by $(\mathcal{P}\qstate)[\selvariable]$ which has the coordinates
\begin{align*}
    \left(\mathcal{P}_{\subspaceof{}}\qstate\right)\left[\selvariable=0\right]
    &= \contraction{\left(\qstatewith\right)^{\dagger},\qstateofat{\goodstatesymbol}{\shortcatvariables}}\\
    \left(\mathcal{P}_{\subspaceof{}}\qstate\right)\left[\selvariable=1\right]
    &= \contraction{\left(\qstatewith\right)^{\dagger},\qstateofat{\goodstatesymbol}{\shortcatvariables}} \, .
\end{align*}
We define a matrix representation of the Grover operator as
\begin{align*}
    \groverofat{\mathrm{red}}{\cselvariable{\outsymbol},\cselvariable{\insymbol}}
    = \coloredmatrixof{
        \cosof{\rotanglesymbol} & -\sinof{\rotanglesymbol} \\
        \sinof{\rotanglesymbol} & \cosof{\rotanglesymbol}
    }{\cselvariable{\outsymbol},\cselvariable{\insymbol}}
\end{align*}
and it holds for any quantum state $\qstatewith$ that
\begin{align*}
    \left(\mathcal{P}_{\subspaceof{}}\groversymbol\qstate\right)\left[\cselvariable{\outsymbol}\right]
    = \contractionof{\groverofat{\mathrm{red}}{\cselvariable{\outsymbol},\cselvariable{\insymbol}},
        \left(\mathcal{P}_{\subspaceof{}}\qstate\right)\left[\cselvariable{\insymbol}\right]
    }{\outsymbol} \, .
\end{align*}
The Grover operator therefore has the eigenvalues:
\begin{itemize}
    \item $\expof{\pm i \rotanglesymbol}$ with the eigen spaces spanned by $\qstateof{\goodstatesymbol}\pm i \qstateof{\badstatesymbol}$
    \item $1$ with the eigen space by the orthogonal complement of $\subspaceof{}$
\end{itemize}
Application the Grover operator $\repnum\in\nn$ times on $\qstateof{0}$ results in the state
\begin{align*}
    \qstateofat{\repnum}{\shortcatvariables}
    = \sinof{\left(\frac{1}{2}+\repnum\right)\rotanglesymbol} \qstateofat{\goodstatesymbol}{\shortcatvariables}
    + \cosof{\left(\frac{1}{2}+\repnum\right)\rotanglesymbol} \qstateofat{\badstatesymbol}{\shortcatvariables} \, .
\end{align*}

\begin{figure}[t]
    \begin{center}
        \begin{tikzpicture}[>=stealth, scale=4]

    % Coordinates
    \coordinate (O) at (0,0);
    \coordinate (W) at (0,1);     % Target state |w>
    \coordinate (Sprime) at (1,0); % All other states |s'>

    % Draw Axes
    \draw[->, thick] (-0.1,0) -- (1.1,0) node[right] {\corelabelsize $\qstateof{\badstatesymbol}$};
    \draw[->, thick] (0,-0.1) -- (0,1.1) node[above] {\corelabelsize $\qstateof{\goodstatesymbol}$};

    \def\sthalfangle{15}
    \def\nextangle{3*\sthalfangle} % 3 * theta
    \draw[->, thick] (O) -- (\sthalfangle:0.9) node[right] {\corelabelsize $\qstateof{0}=\sinof{\frac{\rotanglesymbol}{2}}\qstateof{\goodstatesymbol}+\cosof{\frac{\rotanglesymbol}{2}}\qstateof{\badstatesymbol}$};
    \draw[dashed, thick] (0.4,0) arc (0:\sthalfangle:0.4) node[midway, right, xshift=2pt] {\corelabelsize $\frac{\rotanglesymbol}{2}$};

    \draw[<-, thick, blue] (-1*\sthalfangle:0.6) arc (-\sthalfangle:\sthalfangle:0.6) node[midway, right, yshift=-5pt, xshift=2pt] {\corelabelsize $\rotanglesymbol$};

    % Draw the Oracle Reflection (reflected across s' axis)
    \draw[->, thick, blue] (O) -- (-\sthalfangle:0.9) node[right] {\corelabelsize $-\sinof{\frac{\rotanglesymbol}{2}}\qstateof{\goodstatesymbol}+\cosof{\frac{\rotanglesymbol}{2}}\qstateof{\badstatesymbol}$};


    \draw[->, thick, red] (-1*\sthalfangle:0.8) arc (-\sthalfangle:3*\sthalfangle:0.8) node[midway, right, yshift=15pt, xshift=-5pt] {\corelabelsize $2\rotanglesymbol$};

    % Draw the first Grover Iteration (rotated by 2*theta from original |s>)

    \draw[->, red, thick] (O) -- (\nextangle:0.9) node[above right] {\corelabelsize $\qstateof{1}=\sinof{\left(\frac{1}{2}+1\right)\rotanglesymbol}\qstateof{\goodstatesymbol}+\cosof{\left(\frac{1}{2}+1\right)\rotanglesymbol}\qstateof{\badstatesymbol}$};

    \node[below left] at (O) {\corelabelsize $0$};

\end{tikzpicture}
    \end{center}
    \caption{Grover operator as a rotation in $\spanof{\qstateof{\goodstatesymbol},\qstateof{\badstatesymbol}}$.}
    \label{fig:rotationSketch}
\end{figure}

\subsubsection{Angle measurement}\glsadd{qc:QPE:grover}

Using \gls{qc:QPE} we estimate the angle of the Grover rotation.

\subsubsection{Circuit Realization of the reflection operators}

Let us now investigate how the reflection operations can be realized by circuits.


% THIS WOULD MAKE A DIFFERENCE IN QUANTUM PHASE ESTIMATION!

%\subsubsection{Construction based on controlled Pauli-Z}

First, the reflection about the state $\onehotmapofat{\onetuple{[\catorder]}}{\shortcatvariables}$ is realized by a multiple controlled Pauli-Z (see \figref{fig:reflectionZRealization})
%    \begin{align*}
%        -\reflectionofat{\goodstatesymbol}{\inshortcatvariables,\outshortcatvariables}
%        = \mathrm{MC}^{\seldim-1}\paulizat{\cavariableof{\seldim-1}{\insymbol},\cavariableof{\seldim-1}{\outsymbol},\avariableof{[\seldim-1]}}
%    \end{align*}
\begin{align*}
    -\reflectionofat{\onehotmapof{\onetuple{[\catorder]}}}{\inshortcatvariables,\outshortcatvariables}
    = \contractionof{
        \mathrm{MC^{\catorder-1}}\paulizat{\ccatvariableof{\catorder-1}{\insymbol},\ccatvariableof{\catorder-1}{\outsymbol},\catvariableof{[\catorder-1]}},
        \identityat{\ccatvariableof{[\catorder-1]},\ccatvariableof{[\catorder-1]}{\insymbol},\ccatvariableof{[\catorder-1]}{\outsymbol}}
    }{\inshortcatvariables,\outshortcatvariables}
\end{align*}
Note that a global phase of $-1$ is introduced, which leaves the measurement distribution invariant.
% Extending to e_1
Adding layers of Pauli-X gates to each $\catenumeratorin$ before and after $-\reflectionofat{\onehotmapof{\onetuple{[\catorder]}}}{\inshortcatvariables,\outshortcatvariables}$ (see \figref{fig:reflectionZRealization}) prepares the negated reflection $-\reflectionofat{\onehotmapof{\zerotuple{[\catorder]}}}{\inshortcatvariables,\outshortcatvariables}$, since
\begin{align*}
    \tbasisat{\ccatvariableof{\catenumerator}{\outsymbol}}
    = \contractionof{
        \paulizat{\ccatvariableof{\catenumerator}{\outsymbol},\ccatvariableof{\catenumerator}{\insymbol}},\fbasisat{\ccatvariableof{\catenumerator}{\insymbol}}
    }{\ccatvariableof{\catenumerator}{\outsymbol}} \, .
\end{align*}
% Extending to U
Having a unitary $\unitarysymbol$ preparing $\qstate$ by action on the ground state, we realize the reflection (see \figref{fig:reflectionZRealization})
\begin{align*}
    \reflectionofat{\qstate}{\inshortcatvariables,\outshortcatvariables}
    = \contractionof{
        \unitaryofat{\dagger}{\inshortcatvariables,\ccatvariableof{[\catorder]}{\midsymbol,0}},
        \reflectionofat{\onehotmapof{\zerotuple{[\catorder]}}}{\ccatvariableof{[\catorder]}{\midsymbol,0},\ccatvariableof{[\catorder]}{\midsymbol,1}},
        \unitaryat{\ccatvariableof{[\catorder]}{\midsymbol,1},\outshortcatvariables}
    }{\inshortcatvariables,\outshortcatvariables} \, .
\end{align*}


\begin{figure}[t]
    \begin{center}
        \begin{tikzpicture}[scale=0.35, thick] % , baseline = -3.5pt

%    \drawsmalltensor{0}{9}{\corelabelsize $\fbasis$}
%    \draw (1,9) -- (2,9);
%    \drawsmalltensor{3}{9}{\corelabelsize $\paulixsymbol$}
%    \draw (4,9) -- (5,9);
%    \drawsmalltensor{6}{9}{\corelabelsize $\hgate$}
%
%    %\draw (-1,5) rectangle (1,7);
%    \drawsmalltensor{0}{6}{\corelabelsize $\fbasis$}
%    \draw (1,6) -- (2,6);
%    \drawsmalltensor{0}{2}{\corelabelsize $\fbasis$}
%    \draw (1,2) -- (2,2);
%    \draw (2,1) rectangle (4,7);
%    \drawtextnode{3}{4}{\corelabelsize $\unitarysymbol$}
%    \draw (4,6) -- (12,6);
%    \draw (4,2) -- (12,2);
%
%    %\draw (2,1) rectangle (4,3);
%    %\drawtextnode{3}{2}{\corelabelsize $\hgate$}
%
%
%    \drawtextnode{-3}{9}{\corelabelsize $\headvariable$}
%    \draw[dashed] (-2,7.5) -- (30.5,7.5);
%    \drawtextnode{-3}{4}{\corelabelsize $\shortcatvariables$}
%    \drawtextnode{0}{4.5}{\corelabelsize $\vdots$}
%
%    \draw[dashed] (7.5,12) -- (7.5,-1) node[below] {\corelabelsize $\qstateofat{0}{\shortcatvariables} \otimes \hminusbasat{\headvariable}$};
%
%    \begin{scope}[shift={(4,0)}]
%        \draw (3,9) -- (4,9);
%        \draw (4,8) rectangle (7,10);
%        \node[anchor=center] (text) at (5.5,9) {\corelabelsize $\qcbencodingof{\exformula}$};
%        \draw (5,6) -- (5,8);
%        \drawvariabledot{5}{6}
%        \draw (6,2) -- (6,8);
%        \drawvariabledot{6}{2}
%    \end{scope}
%
%    \drawtextnode{9.5}{11}{\corelabelsize $\reflectionof{\badstatesymbol}$}
%
%    \draw[dashed] (11.5,11) -- (11.5,0.5);

    \drawtextnode{9.5}{4}{\corelabelsize $\inshortcatvariables$}
    \draw (10,6) -- (12,6);
    \drawtextnode{11.5}{4.25}{$\vdots$}
    \draw (10,2) -- (12,2);
    \draw (12,1) rectangle (14,7);
    \drawtextnode{13}{4}{\corelabelsize $\unitaryof{\dagger}$}
    \draw (14,6) -- (15,6);
    \drawsmalltensor{16}{6}{\corelabelsize $\paulixsymbol$}
    \draw (14,2) -- (15,2);
    \drawsmalltensor{16}{2}{\corelabelsize $\paulixsymbol$}
    \draw (17,2) -- (21,2);
    \draw (17,6) -- (18,6);
    \drawsmalltensor{19}{6}{\corelabelsize $\paulizsymbol$}
    \draw (18.5,5) -- (18.5,4);
    \drawvariabledot{18.5}{4}
    \node[rotate=-60] at (18.9,2.75) {$\cdots$};
    \draw (19.5,5) -- (19.5,2);
    \drawvariabledot{19.5}{2}

    \draw (20,6) -- (21,6);
    \drawsmalltensor{22}{6}{\corelabelsize $\paulixsymbol$}
    \draw (23,6) -- (24,6);
    \drawsmalltensor{22}{2}{\corelabelsize $\paulixsymbol$}
    \draw (23,2) -- (24,2);
    \draw (24,1) rectangle (26,7);
    \drawtextnode{25}{4}{\corelabelsize $\unitarysymbol$}
    \draw (26,6) -- (28,6);
    \draw (26,2) -- (28,2);
    \drawtextnode{27}{4.25}{$\vdots$}
    \drawtextnode{28.5}{4}{\corelabelsize $\outshortcatvariables$}
    %\draw (11,9)--(28,9);

%    \drawtextnode{19}{11}{\corelabelsize $\reflectionof{0}$}

    \draw[dashed] (23.5,0) -- (23.5,9.25) -- (14.5,9.25) -- (14.5,0) -- cycle;
    \drawtextnode{16}{8.5}{\corelabelsize $-\reflectionof{\onehotmapof{\zerotuple{[\catorder]}}}$}

    \draw[dashed] (20.5,1.25) -- (20.5,8.75) -- (17.5,8.75) -- (17.5,1.25) -- cycle;
    \drawtextnode{19}{8}{\corelabelsize $-\reflectionof{\onehotmapof{\onetuple{[\catorder]}}}$}

    \begin{scope}[shift={(-12,0)}]
        \drawtextnode{9.5}{4}{\corelabelsize $\inshortcatvariables$}
        \draw (10,6) -- (12,6);
        \drawtextnode{11.5}{4.25}{$\vdots$}
        \draw (10,2) -- (12,2);
        \draw (12,1) rectangle (14,7);
        \drawtextnode{13}{4}{\corelabelsize $-\reflectionof{\qstate}$}
        \draw (14,6) -- (15,6);

        \draw (14,6) -- (16,6);
        \drawtextnode{14.5}{4.25}{$\vdots$}
        \draw (14,2) -- (16,2);
        \drawtextnode{16.5}{4}{\corelabelsize $\outshortcatvariables$}

        \drawtextnode{19}{4}{${=}$}
    \end{scope}


    %\draw[dashed]  -- ;
    %\draw[dashed]  -- ; % node[below] {\corelabelsize $\qstateofat{1}{\shortcatvariables} \otimes \hminusbasat{\headvariable}$};


\end{tikzpicture}
    \end{center}
    \caption{
        Preparation of the negated reflection operator $-\reflectionof{\qstate}$ about a state $\qstate$, which is prepared by a unitary $\unitarysymbol$ acting on the ground state $\onehotmapofat{\zerotuple{[\catorder]}}{\shortcatvariables}$.
        The multiply controlled Pauli-Z is the negated reflection $-\reflectionof{\onetuple{[\catorder]}}$, and extended to the negated reflection $-\reflectionof{\onetuple{[\catorder]}}$ about the ground state by adding layers of single qubit Pauli-X.
        When concatenating with $\unitaryof{\dagger}$ and $\unitarysymbol$, the $-\reflectionof{\qstate}$ results.
    }
    \label{fig:reflectionZRealization}
\end{figure}

%\subsubsection{Construction based on oracle with ancilla head qubit}

Alternatively we can use the phasekick technique when acting on an disentangled ancilla qubit prepared in the Hadamard basis state
\begin{align*}
    \hminusbasat{\headvariable}
    = \frac{1}{\sqrt{2}} \cdot \coloredmatrixof{1 \\ -1}{\headvariable} \, .
\end{align*}
We notice, that
\begin{align*}
    \contractionof{
        \qcbencodingofat{\exformula}{\cheadvariable{\outsymbol},\cheadvariable{\insymbol},\indexedshortcatvariables},
        \hminusbasat{\cheadvariable{\insymbol}}
    }{\cheadvariable{\outsymbol}}
    = \begin{cases}
          \hminusbasat{\cheadvariable{\insymbol}} & \ifspace \exformula(\shortcatindices) = 0\\
          -\hminusbasat{\cheadvariable{\insymbol}} & \ifspace \exformula(\shortcatindices) = 1
    \end{cases}
\end{align*}

Let us consider the span of one-hot encoded models of $\lnot\exformula$, namely
\begin{align*}
    \subspaceof{\lnot\exformula}
    \coloneqq \spanof{\onehotmapofat{\shortcatindices}{\shortcatvariables}\wcols\exformula(\shortcatindices)=0}
\end{align*}
and the span of one-hot encoded models
\begin{align*}
    \subspaceof{\exformula}
    \coloneqq\spanof{\onehotmapofat{\shortcatindices}{\shortcatvariables}\wcols\exformula(\shortcatindices)=1} \, .
\end{align*}
The reflection across $\subspaceof{\lnot\exformula}$ is then performed by (see \figref{fig:reflectionCCRealization})
\begin{align*}
    \reflectionofat{\subspaceof{\lnot\exformula}}{\outshortcatvariables,\inshortcatvariables}
    = \contractionof{
        \hminusbasat{\cheadvariable{\outsymbol}},
        \qcbencodingofat{\exformula}{\cheadvariable{\outsymbol},\cheadvariable{\insymbol},\shortcatvariables},
        \hminusbasat{\cheadvariable{\insymbol}},
        \identityat{\shortcatvariables,\outshortcatvariables,\inshortcatvariables}
    }{\outshortcatvariables,\inshortcatvariables}
\end{align*}
We further have
\begin{align*}
    \reflectionofat{\subspaceof{\exformula}}{\outshortcatvariables,\inshortcatvariables}
    = -\reflectionofat{\left(\subspaceof{\exformula}\right)^{\perp}}{\outshortcatvariables,\inshortcatvariables}
    = -\reflectionofat{\subspaceof{\lnot\exformula}}{\outshortcatvariables,\inshortcatvariables} \, .
\end{align*}
When preparing a disentangled ancilla qubit $\headvariable$ in the basis state $\hminusbasat{\headvariable}$, the oracle $\qcbencodingof{\exformula}$ is effectively a negated reflection across the span of one-hot encoded models of $\exformula$.

%This is an effective
%This further coincides with the effective negative reflection across the span of one-hot encoded models
%\begin{align*}
%    \subspaceof{\exformula}
%    \coloneqq\spanof{\onehotmapofat{\shortcatindices}{\shortcatvariables}\wcols\exformula(\shortcatindices)=1} \, .
%\end{align*}

%We will investigate different effective schemes to realize $\reflectionof{\badstatesymbol}$ later.
%\begin{itemize}
%    \item When $\qstateof{\goodstatesymbol}=\onehotmapofat{\onetuple{[\seldim]}}{\avariables}\otimes $ and $\qstate$
%    \item Reflection across the hyperplane $\spanof{\onehotmapofat{\shortcatindices}{\shortcatvariables}\wcols\exformula(\shortcatindices)=0}$ effectively by $\qcbencodingof{\exformula}$ acting on an auxiliary disentangled qubit in $\hminusbasat{\headvariable}$
%\end{itemize}

\begin{figure}[t]
    \begin{center}
        \begin{tikzpicture}[scale=0.35, thick] % , baseline = -3.5pt


    %\drawtextnode{9.5}{4}{\corelabelsize $\inshortcatvariables$}

    \drawsmalltensor{10}{8}{\corelabelsize $\hminusbas$}
    \draw (11,8) -- (12,8);
    \draw (12,7) rectangle (15,9);
    \drawtextnode{13.5}{8}{\corelabelsize $\qcbencodingof{\formula}$}
    \draw (13,7) -- (13,6);
    \drawvariabledot{13}{6}
    \draw (14,7) -- (14,2);
    \drawvariabledot{14}{2}
    \draw (15,8) -- (16,8);
    \drawsmalltensor{17}{8}{\corelabelsize $\hminusbas$}

    \drawtextnode{10.5}{4}{\corelabelsize $\inshortcatvariables$}
    \draw (11,6) -- (16,6);
    \drawtextnode{12.5}{4.25}{$\vdots$}
    \draw (11,2) -- (16,2);
    \drawtextnode{16.5}{4}{\corelabelsize $\outshortcatvariables$}
    %\draw (12,1) rectangle (14,7);
    %\drawtextnode{13}{4}{\corelabelsize $\unitaryof{\dagger}$}
    %\draw (14,6) -- (15,6);

    \begin{scope}[shift={(-12,0)}]
        \drawtextnode{9.5}{4}{\corelabelsize $\inshortcatvariables$}
        \draw (10,6) -- (12,6);
        \drawtextnode{11.5}{4.25}{$\vdots$}
        \draw (10,2) -- (12,2);
        \draw (12,1) rectangle (14,7);
        \drawtextnode{13}{4}{\corelabelsize $-\reflectionof{\subspaceof{\formula}}$}
        \draw (14,6) -- (15,6);

        \draw (14,6) -- (16,6);
        \drawtextnode{14.5}{4.25}{$\vdots$}
        \draw (14,2) -- (16,2);
        \drawtextnode{16.5}{4}{\corelabelsize $\outshortcatvariables$}

        \drawtextnode{19}{4}{${=}$}
    \end{scope}


    %\draw[dashed]  -- ;
    %\draw[dashed]  -- ; % node[below] {\corelabelsize $\qstateofat{1}{\shortcatvariables} \otimes \hminusbasat{\headvariable}$};


\end{tikzpicture}
    \end{center}
    \caption{
        Realization of the reflection across the subspace of one-hot encoded models to a Boolean function $\formula$.
        An ancilla qubit is prepared in the Hadamard basis state $\hminusbas$ and serves as the head qubit of the oracle $\qcbencodingof{\formula}$.
        By the phase kickback effect, the qubit stays disentangled and the reflection $-\reflectionof{\subspaceof{\formula}}$ is effectively performed on the qubits $\shortcatvariables$.
    }
    \label{fig:reflectionCCRealization}
\end{figure}
